Across the sixteen languages analyzed, orthographic similarity ranges from
0.3423 to 0.9391, while phonetic similarity occupies a much narrower and higher
band (0.7622 to 0.9927). This discrepancy reveals a fundamental dynamic in
Philippine linguistics: spelling has diverged far more rapidly than sound.
Orthography is highly susceptible to external standardization, bearing the
marks of disparate missionary, educational, and administrative decisions even
among neighboring communities. Phonology, by contrast, is markedly more
conservative. The spoken word preserves a shared Austronesian history that
modern spelling conventions obscure. Read against the linguistic map
from~\cref{fig:lingmap_out}, it becomes evident that while
geography—specifically rugged terrain—explains part of this pattern, historical
maritime contact and creolization frequently override geographic boundaries.

\subsection{Visayas and Archipelagic Connectivity}
The Bisayan languages form the most cohesive group in the dataset across both
measures. Rather than acting as barriers, the narrow straits of the Visayas
have functioned as conduits for linguistic exchange. This is most pronounced
between Hiligaynon and Cebuano, which exhibit near-total phonetic convergence
(0.9927) alongside high orthographic similarity (0.9391). Kinaray-a and
Waray-Waray remain closely tied to this core; remarkably, Waray-Waray maintains
a 0.9550 phonetic similarity with Kinaray-a despite being geographically
separated by the Samar Sea.

Romblomanon exemplifies this maritime connectivity. Geographically situated
between the northern and southern migration routes, its linguistic profile
firmly anchors it at the center of the Bisayan cluster rather than its
periphery. It shares high phonetic similarity with Hiligaynon (0.9734), Cebuano
(0.9708), and Tagalog (0.9716), reflecting a historical zone where mutual
intelligibility aligns with the ease of sea travel. Similarly, Bikol and
Masbatenyo form a tight pair (0.9206 orthographically, 0.9857 phonetically),
reflecting their adjacency on the Bicol Peninsula.

Further north, Tagalog pairs closely with Cebuano and Hiligaynon in spelling
(reaching 0.8877 and 0.8811, respectively). This orthographic alignment is
likely sustained by modern educational infrastructure and mass media
standardizing the archipelago's most widely spoken languages, rather than deep
structural convergence.

\subsection{Luzon and the Cordillera Barrier}
In the northern Philippines, physical terrain plays a far more restrictive
role, though the linguistic landscape remains complex. Ilocano demonstrates low
similarity to the central languages in both sound (ranging from 0.76 to 0.79)
and spelling (0.55 to 0.56). This stark division tracks closely with the
Cordillera and Sierra Madre mountain ranges, which have historically isolated
Ilocano speakers from central lowland influences.

However, terrain does not account for all northern variations. Paranan and
Pangasinan display an unexpectedly high phonetic affinity (0.9677) despite
geographic separation, suggesting either the retention of conservative features
lost in central varieties, or parallel trajectories of simplification in
isolated communities. Similarly, Kapampangan proves that geographic adjacency
does not guarantee orthographic convergence. Despite a long literary history in
Central Luzon, its spelling diverges significantly from neighboring Ilocano
(0.5421). Phonetically, however, Kapampangan is closer to the distant Kinaray-a
(0.9456) than to its immediate neighbors. This mismatch characterizes lowland
regions, where historical trade, migration, and urbanization blur the strict
linguistic borders imposed by terrain.

\subsection{The Impact of Creolization}
Chavacano emerges as the clearest outlier in the matrix, demonstrating that
creolization reshapes a language far more aggressively than either distance or
ambient contact. As a Spanish-based creole, its orthographic similarity with
neighboring languages largely bottoms out between 0.35 and 0.45. Phonetically,
it sits near the floor of the dataset, scoring as low as 0.7881 against
Tagalog.

Although Chavacano absorbed extensive vocabulary from Cebuano and Tausug
substrates in Zamboanga, its phonotactics and spelling conventions were
fundamentally reconfigured by Spanish colonizers. The limited orthographic
similarity it does share with other Philippine languages is an artifact of the
shared Latin script, not common ancestry. Its profile firmly suggests that the
mechanics of creolization overwrite inherited Austronesian traits.

\subsection{Summary}
Taken together, these measures demonstrate that phonetic and orthographic data
record two different histories. Orthographic similarity serves as a modern
record of which communities learned to write together under shared
administrative or educational standards. Phonetic similarity offers a deeper,
more conservative record of shared descent and historical contact. Geography
predicts these patterns only where terrain—such as the Cordillera—genuinely
restricts human movement. Where communities could navigate waters, engage in
trade, or experience colonization, human contact readily overwrote the
geographic map. Languages that resist geographic explanation on the phonetic
measure, particularly Yami and Paranan, present clear targets for future
diachronic study.